% claim-ref: pothos#bud-growth-time
% claim-ref: pothos#combine-rooted-cuttings
% claim-ref: pothos#cutting-callus-time
% claim-ref: pothos#established-plant-substrate-recipe
% claim-ref: pothos#failure-cue-blackened-leaf-edges
% claim-ref: pothos#failure-cue-brown-roots
% claim-ref: pothos#failure-cue-no-bud-growth
% claim-ref: pothos#medium-placement
% claim-ref: pothos#prepare-node-cutting
% claim-ref: pothos#primary-method
% claim-ref: pothos#propagation-medium
% claim-ref: pothos#rooting-time
% claim-ref: pothos#transplant-readiness
% claim-ref: pothos#water-placement
% claim-ref: pothos#water-root-photo-time
% claim-ref: pothos#watering-cue
\CompanionHeader{PROPAGATION}{Pothos}{Provisional Epipremnum aureum}{Indoors; specimen identity, cultivar, variegation and household conditions remain unverified.}
\Context{Node-bearing stem cutting / warm conditions. General guidance; observe parent readiness before acting.}
\begin{minipage}[t]{.625\linewidth}\vspace{0pt}
\Step{1}{Prepare the rooting space}{Use vermiculite for the primary medium method; this is temporary rooting material, not the established-plant recipe. Water rooting is the alternate below. \textbf{[POT-WISC]}}
\Step{2}{Select healthy material}{Choose a healthy vine section carrying a leaf, a node and an axillary bud. A detached leaf alone is not the specified propagule. \textbf{[POT-NCSU-PROP; POT-WISC]}}
\Step{3}{Locate the node \& cut}{Cut a stem section that retains the node and bud. Track which way the bud faces before placing the section. No fixed cutting length is established here. \textbf{[POT-NCSU-PROP]}}
\Step{4}{Orient the cutting}{Place the node in moist propagation medium with the bud upward and covered; leave the leaf exposed. The schematic shows orientation, not a measured depth. \textbf{[POT-NCSU-PROP]}}
\Step{5}{Maintain the rooting environment}{Keep the vermiculite moist and aerated in a warm rooting location. Observe the leaf and stem condition. The sources do not prescribe a separate callusing interval. \textbf{[POT-NCSU-PROP; POT-WISC]}}
\Step{6}{Separate bud growth from roots}{In warm conditions buds can start in 1-2 weeks; vines root in water or vermiculite in 3-4 weeks. Bud growth alone does not prove transplant readiness. \textbf{[POT-WISC]}}
\Step{7}{Check readiness}{Proposed: pot only after several roots have formed and the cutting sustains growth. There is no published root-length threshold here; inspect rather than transplanting by calendar alone. \textit{Proposed starting point.} \textbf{[POT-NCSU-PROP; POT-WISC]}}
\Step{8}{Transfer to normal care}{Several rooted cuttings may share a container for a fuller pot. Use the established-plant draining mix as a proposed starting point; then follow surface-dry watering cues and observe recovery. \textit{Proposed starting point.} \textbf{[POT-NCSU; POT-WISC]}}
\end{minipage}\hfill\begin{minipage}[t]{.345\linewidth}\vspace{0pt}
\Note{Milestones, not deadlines}{Buds: 1-2 weeks. Roots: 3-4 weeks in warm conditions. NCSU's two-week photograph is an observation, not a guaranteed deadline. \textbf{[POT-NCSU; POT-WISC]}}
\Note{Water alternate}{Root the node in water. Proposed: keep foliage above the water, then transfer after several roots and sustained growth. Water rooting does not establish a permanent hydroponic recipe. \textit{Proposed starting point.} \textbf{[POT-NCSU-PROP; POT-WISC]}}
\Note{Check 1}{Brown/absent roots: check excess water and correct moisture. \textit{Proposed starting point.} \textbf{[POT-WISC]}}
\Note{Check 2}{No bud growth: check warmth; do not discard by the calendar alone. \textit{Proposed starting point.} \textbf{[POT-WISC]}}
\Note{Check 3}{Black margins/tips: check moisture extremes and fertilizer salts. \textit{Proposed starting point.} \textbf{[POT-WISC]}}
\end{minipage}\par\vspace{4pt}
\Panel{Ready for normal care}{Several roots plus sustained growth support transfer. Log the method and conditions; seed timing and a universal clonal lifespan remain unknown.\par Start date: \hrulefill\quad Method: \hrulefill\par First successful growth / conditions: \hrulefill}
\Sources{Evidence: POT-NCSU NC State; POT-NCSU-PROP NC State; POT-WISC University of Wisconsin--Madison Extension. Full titles, URLs and claim scopes: entry sources.yaml and companion worksheets. Conversions are rounded; source units remain authoritative.}{pothos / propagation}
